\chapteropen
\chapter{Title of published work, or similar}\label{ch:paper-example}
\hypertarget{ch:paper-example}{}
\glsresetall
\quoteThree

% Check current Graduate College and publisher requirements before reproducing
% published material.
%
% Depending on the circumstances, this opening material may also need to state:
% - the publication status of the work;
% - whether the chapter is reproduced or adapted;
% - the dissertation author's contribution;
% - the contributions of co-authors;
% - reuse permission or licensing information;
% - changes made from the published version.
% 
% The \chapterpublicationnote will format this consistently

\chapterpublicationnote
    {Pomeroy_2026a} % this should be a normal citation reference to a bibliography item
    {This chapter is adapted from the published article.}
    {This article was published under the Creative Commons Attribution 4.0 International (CC BY 4.0) licence:
    \url{https://creativecommons.org/licenses/by/4.0/}.}
    {The candidate led the analysis, interpretation, figure preparation, and manuscript drafting; co-authors contributed to scientific discussion, supervision, and revision.}


\section{Chapter Introduction}
For a chapter where a published work is included, it is good to have a summary of previous chapters, and what motivates the work in this paper. The paper can be dropped in verbatim (or adapted with formatting changes as appropriate).


\clearpage
\section{Manuscipt: Include the Actual Title from the Published Paper}\label{sec:example_paper}

\begin{center}
    \textbf{Abstract}
\end{center}
% \section*{}
The paper's abstract will sit here

\par
\subsection{Introduction}\label{sec:example_paper_intro}
Now this is where the actual paper should be dropped. Note that the sections in the original paper will be at a lower level, \ie the published work \texttt{section} will become a \texttt{subsection} here, etc.


\clearpage
\section{Chapter Summary}
Finally, for a published paper chapter, it is worthwhile providing a transition paragraph, summarizing the finds, and then outlining the motivation for the next steps of the analysis in the subsequent chapter(s).
